\section{Conclusion}
\label{sec:conclusion}
This study analyzed citation patterns in GESs across political domains, revealing significant differences between GES models and political parties in primary information source usage and highlighting the poisoning attack surface of GESs.
Our method identified publisher attributes and quantified how publisher categories influenced answer generation by examining citation preferences and content reflection faithfulness in GES outputs.
Our results showed that primary information sources comprised \(60\%\)--\(65\%\) in Japan versus \(25\%\)--\(45\%\) in the U.S., and low-barrier sources comprised approximately \(30\%\) of citations, indicating higher poisoning risk for U.S. political answers.
% We further found that low content-injection-barrier sources had reduced faithfulness of reflection in answer content.
% We believe that our study contributes to the development of more secure and reliable GES systems as information delivery channels to citizens.
